\documentclass[12pt]{amsart} 
\usepackage{mathtools}

\input{./latex-preamble/cm-preamble.tex}
\input{./latex-preamble/use-scr.tex}
\input{./latex-preamble/fr-mod.tex}


%\graphicspath{ {./diagrams/} }

\begin{document}

{Hartshorne} \hfill {11-2-2026}

\bigskip\bigskip

\begin{center}

{\sc Chapter IV.3}

\end{center}

\begin{center}

  {Noah Parker} 
    
\end{center}

\bigskip\bigskip

\begin{enumerate}
	\item[1.] {
      If $X$ is a curve of genus 2, show that a divisor is very ample $\iff$ $\deg D\geqs 5$.
      This strengthens (3.3.4).
    }
\end{enumerate} 

\begin{proof}
\end{proof}

\begin{enumerate}
	\item[2.] {
      Let $X$ be a plane curve  of degree 4.
      \begin{enumerate}
        \item Show that the effective canonical divisors on $X$ are exactly the divisors $X.L$, where $L$ is a line in $\PP^2$.
        \item If $D$ is any effective divisor of degree 2 on $X$, show that $\dim|D|=0$.
        \item Conclude that $X$ is not hyperelliptic (Ex. 1.7).
      \end{enumerate}
    }
\end{enumerate}

\begin{proof}
\end{proof}

\begin{enumerate}
	\item[3.] {
      If $X$ is a curve of genus $\geqs 2$ which is a complete intersection (II, Ex. 8.4) in some $\PP^n$, show that the canonical divisor $K$ is very ample.
      Conclude that a curve of genus 2 can never be a complete intersection in any $\PP^n$. Cf. (Ex 5.1).
    }
\end{enumerate}

\begin{proof}
\end{proof}

\begin{enumerate}
	\item[4.] {
      Let $X$ be the $d$-uple embedding (I, Ex. 2.12) of $\PP^1$ in $\PP^d$, for any $d\geqs 1$.
      We call $X$ the \emph{rational normal curve of degree $d$} in $\PP^d$.
      \begin{enumerate}
        \item Show that $X$ is projectively normal, and that its homogeneous ideal can be generated by forms of degree 2.
      \end{enumerate}
    }
\end{enumerate}

\begin{proof}
  Recall that $X\subset \PP^r_A$ is \emph{projectively normal} for the given embedding, if its homogeneous coordinate ring $S(X)$ is an integrally closed domain.
\end{proof}

\begin{enumerate}
	\item[5.] {
      Let $X$ be a curve in $\PP^3$ which is not contained in any plane.
      \begin{enumerate}
        \item If $O\notin X$ is a point, such that the projection from $O$ induces a birational morphism $\vphi$ from $X$ to its image in $\PP^2$, show that $\vphi(X)$ \emph{must} be singular.
          [\emph{Hint}: Calculate $\dim H^0(X,\O_X(1))]$ in two ways.]
        \item If $X$ has degree $d$ and genus $g$, conclude that $g<\frac{1}2(d-1)(d-2)$. (Use (Ex. 1.8).)
        \item Now let $\{X_t\}$ be the flat famiily of curves induced by the projection (III, 9.8.3) whose fibre over $t=1$ is $X$, and whose fibre $X_0$ over $t=0$ is a scheme with support $\vphi(X)$.
          Show that $X_0$ always has nilpotent elements.
          Thus the example (III, 9.8.4) is typical.
      \end{enumerate}
    }
\end{enumerate} 

\begin{proof}
\end{proof}

\begin{enumerate}
	\item[11a.] {
      If $X$ is a nonsingular variety of dimension $r$ in $\PP^n$, and if $n>2r+1$, show there is a point $O\notin X$ such that the projection from $O$ induces a closed immersion of $X$ into $\PP^{n-1}$.
    }
\end{enumerate}

\begin{proof}
\end{proof}

\end{document}
